\section{Item LR(0), Closure, dan Goto}
\label{sec:ch12-item}

Konstruksi item, closure, dan goto pada bagian ini mengikuti Mogensen
Bab 3.14 dan Thain Bab 4.4.

\subsection{Item LR(0)}

\begin{definisi}[Item LR(0)]
Item LR(0) adalah aturan produksi dengan sebuah titik yang menunjukkan
posisi parser. Item \([A \to \alpha \cdot \beta]\) berarti parser
telah mengenali \(\alpha\) dan mengharapkan \(\beta\).
\end{definisi}

Sebagai contoh, aturan \(E \to E + T\) memiliki empat item:
\[
[E \to \cdot E + T],\quad
[E \to E \cdot + T],\quad
[E \to E + \cdot T],\quad
[E \to E + T \cdot].
\]
Titik di awal berarti belum ada yang dikenali; titik di akhir berarti
seluruh ruas kanan sudah dikenali dan aturan siap direduksi.

\subsection{Closure}

Jika item \([A \to \alpha \cdot B \beta]\) berada dalam himpunan
item dan \(B\) adalah nonterminal, maka semua aturan \(B \to \gamma\)
harus ditambahkan dalam bentuk \([B \to \cdot \gamma]\). Proses ini
disebut closure.

\begin{contoh}[Closure]
Untuk item \([E \to \cdot T]\) dan aturan
\(T \to T * F \mid F\), closure menambahkan
\([T \to \cdot T * F]\) dan \([T \to \cdot F]\). Karena \(F\)
juga nonterminal, closure menambahkan aturan \(F\) pula.
\end{contoh}

\subsection{Goto}

\(\text{goto}(I, X)\) adalah closure dari himpunan item yang titiknya
bergeser melewati simbol \(X\). Jika \(I\) memuat
\([A \to \alpha \cdot X \beta]\), maka
\(\text{goto}(I,X)\) memuat \([A \to \alpha X \cdot \beta]\).

\begin{figure}[H]
\centering
\begin{tikzpicture}[node distance=6mm and 10mm]
  \node[phase, text width=5.0cm] (i0) {\(I_0\):\\\([E \to \cdot E + T]\)\\
        \([E \to \cdot T]\)};
  \node[phase, text width=5.0cm, right=of i0] (i1) {\(I_1\):\\\([E \to E \cdot + T]\)};
  \draw[arrow] (i0) -- node[above, font=\small] {\(E\)} (i1);
\end{tikzpicture}
\caption{Item dan transisi goto pada simbol \(E\)}
\label{fig:ch12-item}
\end{figure}

\subsection{Koleksi Kanonik}

Koleksi kanonik LR(0) dibangun dengan:
\begin{enumerate}[label=\arabic*.,leftmargin=*]
  \item mulai dari closure item awal \([S' \to \cdot S]\);
  \item untuk setiap state dan simbol \(X\), hitung
        \(\text{goto}(I,X)\);
  \item jika hasilnya belum ada, tambahkan sebagai state baru;
  \item ulangi sampai tidak ada state baru.
\end{enumerate}

Jumlah state bergantung pada grammar. Untuk grammar ekspresi pada
\file{grammar-lr.bnf}, koleksi kanonik memiliki 12 state.
